% Gleichungen lösen · Schnittpunkte und Stellen · Prüfungs-Fokus Abitur GK – bank/gleichungen-loesen/e1.jsonl (zusammenbau v0.6, Prüfungs-Fokus)
\einheitenkopf[e1]{Einheit 1 · Ganzrationale Gleichungen}
\zweigzeile{Abi GK ×4}
\setcounter{aufgabe}{0}

% Anlauf (ohne Prüfkennung)
\begin{aufgabe}{Schnittpunkte und Stellen – Anlauf}
\begin{teile}
\teil Welcher Weg? $x^4 - 10x^2 + 9 = 0$ \kreuz{Lösungsformel} \\ \kreuz{Substitution} \\ \kreuz{Polynomdivision}
\teil $f(x) = x^2 - 2x - 1$, $g(x) = x + 9$. Schnittpunkte? S1( \leerfeld | \leerfeld ), S2( \leerfeld | \leerfeld )
\end{teile}
\end{aufgabe}

\begin{aufgabe}{Schnittpunkte und Stellen – Prüfungsaufgaben}
\begin{teile}
\teil $f(x) = 0{,}5x^3 - 2x + 3$. An welchen Stellen schneidet der Graph die Gerade $y = 3$? \hfill (A19) \\ x1 = \leerfeld , x2 = \leerfeld , x3 = \leerfeld
\teil $f(x) = x^3 + 2x^2 - 8x - 1$. An welchen Stellen schneidet der Graph die Gerade $y = -1$? \hfill (A19) \\ x1 = \leerfeld , x2 = \leerfeld , x3 = \leerfeld
\teil Ein Mountainbiker springt von einer Rampe. Seine Flugbahn ist $f(x) = -0{,}02x^2 + 20$, das Gelände $g(x) = 0{,}0001x^4 - 0{,}05x^2 + 16$ ($x \ge 0$, 1 LE = 1 m). Koordinaten des Landepunkts $L$? \hfill (A18) \\ L( \leerfeld | \leerfeld )
\teil $f(x) = x^2 + 3x - 1$, $g(x) = -x^2 - x + 5$. Zeige, dass sich die Graphen nur für $x = -3$ und $x = 1$ schneiden. \hfill (A19)
\teil $f(x) = 0{,}5x^2 - x + 2$, $g(x) = -0{,}5x^2 + 3x + 2$. Zeige, dass sich die Graphen nur für $x = 0$ und $x = 4$ schneiden. \hfill (A19)
\teil $f(x) = 2x^2 - x - 7$, $g(x) = x^2 + 2x + 3$. Zeige, dass sich die Graphen nur für $x = -2$ und $x = 5$ schneiden. \hfill (A19)
\teil Ein Blatt in einer Grafik wird von $u(x) = 0{,}5x^3$ und $v(x) = x^2 \cdot (3 - x)$ begrenzt. Zeige, dass $P(0 | 0)$ und $Q(2 | 4)$ die einzigen gemeinsamen Punkte sind. \hfill (A20)
\teil Eine Flosse wird von $u(x) = 0{,}25x^3$ und $v(x) = 0{,}5x^2 \cdot (6 - x)$ begrenzt. Zeige, dass $P(0 | 0)$ und $Q(4 | 16)$ die einzigen gemeinsamen Punkte sind. \hfill (A20)
\teil Ein Segel wird von $u(x) = x^3$ und $v(x) = 2x^2 \cdot (3 - x)$ begrenzt. Zeige, dass $P(0 | 0)$ und $Q(2 | 8)$ die einzigen gemeinsamen Punkte sind. \hfill (A20)
\teil $h(x) = \frac{4}{x - 2}$ ($x \ne 2$), $g(x) = x - 2$. Koordinaten der gemeinsamen Punkte? \hfill (A19) \\ S1( \leerfeld | \leerfeld ), S2( \leerfeld | \leerfeld )
\teil $h(x) = \frac{9}{x + 4}$ ($x \ne -4$), $g(x) = x + 4$. Koordinaten der gemeinsamen Punkte? \hfill (A19) \\ S1( \leerfeld | \leerfeld ), S2( \leerfeld | \leerfeld )
\teil $h(x) = \frac{2}{x - 3}$ ($x \ne 3$), $g(x) = 0{,}5 \cdot (x - 3)$. Koordinaten der gemeinsamen Punkte? \hfill (A19) \\ S1( \leerfeld | \leerfeld ), S2( \leerfeld | \leerfeld )
\end{teile}
\end{aufgabe}

\begin{aufgabe}{Schnittpunkte und Stellen – Prüfungsaufgaben – weiter}
\begin{teile}
\teil Zeige: $-0{,}5 \cdot (x - 2) \cdot (x + 3)^2 = -0{,}5x^3 - 2x^2 + 1{,}5x + 9$. \hfill (A20)
\teil Zeige: $0{,}2 \cdot (x + 1) \cdot (x - 3)^2 = 0{,}2x^3 - x^2 + 0{,}6x + 1{,}8$. \hfill (A20)
\teil Zeige: $-2 \cdot (x + 4) \cdot (x - 1)^2 = -2x^3 - 4x^2 + 14x - 8$. \hfill (A20)
\end{teile}
\end{aufgabe}

\medskip
%% Merkkasten Einheit 1: 8 Zeilen, Vorgabe höchstens fünf (3.1) – ungekürzt
\uebersichtskasten{Gleichsetzen: Schnittpunkte zweier Graphen sind die Lösungen von f(x) = g(x); die Stelle zu einem Wert die Lösung von f(x) = c. Alles auf eine Seite bringen und ordnen, dann das Verfahren wählen. Ein Schnittpunkt hat zwei Koordinaten – den Funktionswert nachrechnen. \\ Ausklammern und Nullprodukt: Ohne Absolutglied x (oder x²) ausklammern; ein Produkt ist null, wenn ein Faktor null ist – x = 0 ist dann eine Lösung, nicht durch x teilen. \\ f(x) = x³ − x und g(x) = 3x: x³ − 4x = 0 ⇔ x · (x² − 4) = 0 ⇔ x = 0 oder x = 2 oder x = −2; Schnittpunkte (0 | 0), (2 | 6), (−2 | −6). \\ Lösungsformel: x² + px + q = 0 hat die Lösungen x = −p/2 ± √((p/2)² − q) – höchstens zwei; das begründet ein „nur“. \\ Substitution: Nur gerade Exponenten (biquadratisch): u = x², die quadratische Gleichung in u lösen, nur u ≥ 0 zurücksubstituieren, dann x = ±√u. \\ x⁴ − 5x² + 4 = 0: u² − 5u + 4 = 0 ⇔ u = 1 oder u = 4 ⇔ x = ±1 oder x = ±2. \\ n-te Wurzel: xⁿ = c ⇔ x = ⁿ√c, bei geradem n auch x = −ⁿ√c. \\ Polynomdivision: Ist eine Lösung x₁ bekannt (aus der Aufgabe oder durch Probieren), das Polynom durch (x − x₁) teilen; der Rest ist eine quadratische Gleichung.}
